\annotate

% --- Часть 1: Сведения об общем объёме работы и структуре рукописи ---

Работа содержит AAA~страниц, BBB~книг отчета, CCC~рисунков, DDD~таблиц,
EEE~источников цитирования, FFF~приложений.

% Допустимы сокращения:
% Отчет 85 с., 1 кн., 24 рис., 12 табл., 50 источн., 2 прил.

% --- Часть 2: Перечень ключевых слов ---

ШАБЛОН, LATEX, ОТЧЕТ, КУРСОВОЙ ПРОЕКТ, ВКР

% Перечень ключевых слов должен включать от 5 до 15 слов или словосочетаний из
% текста отчета, которые в наибольшей мере характеризуют его содержание и
% обеспечивают возможность информационного поиска. Они приводятся в
% именительном падеже и печатаются прописными буквами, в строку, через запятые,
% без абзацного отступа и переноса слов, без точки в конце перечня.

% --- Часть 3: Текст реферата ---

Описание предмета, темы, цели работы. Предмет, тема, цель работы указываются в
том случае, если они не ясны из заглавия документа.

Описание метода или методологии проведения работы. Метод или методологию
проведения работы целесообразно описывать в том случае, если они отличаются
новизной или представляют интерес с точки зрения данной работы. Широко
известные методы только называются. В рефератах документов, описывающих
экспериментальные работы, указывают источники данных и характер их обработки.

Описание результатов работы. Результаты работы описывают предельно точно и
информативно. Приводятся основные теоретические и экспериментальные результаты,
фактические данные, обнаруженные взаимосвязи и закономерности. При этом
отдается предпочтение новым результатам и данным долгосрочного значения, важным
открытиям, выводам, которые опровергают существующие теории, а также данным,
которые по мнению автора документа имеют практическое значение. Следует указать
пределы точности и надежности данных, а также степень их обоснования. Уточняют,
являются ли цифровые значения первичными или производными, результатом одного
наблюдения или повторных испытаний.

Описание области применения результатов. Область применения результатов важно
указывать для патентных документов.

Описание выводов. Выводы могут сопровождаться рекомендациями, оценками,
предложениями, гипотезами, описанными в исходном документе.

% Описание рекомендаций по внедрению или итогов внедрения результатов работы.

% Описание экономической эффективности или значимости работы.

% Описание прогнозных предположений о развитии объекта исследования.

Дополнительная информация. Дополнительная информация включает данные, не
существенные для основной цели исследования, но имеющие значение вне его
основной темы. Кроме того, можно указывать название организации, в которой
выполнена работа, сведения об авторе исходного документа, ссылки на ранее
опубликованные документы и т.п.


% Если отчет не содержит сведений по какой-либо из перечисленных структурных
% частей реферата, то в тексте реферата она опускается, при этом
% последовательность изложения сохраняется.

% Оптимальный объем текста реферата - 850 печатных знаков, но не более одной
% страницы машинописного текста.


\newpage
\nchapter{Abstract}

Перевод реферата на английский язык (для специалистов и магистров).
